%% !TEX root = manuscript.tex

\chapter{Activités complémentaires}
\minitoc
\newpage

En parallèle de mes deux grandes missions, mon alternance a compté d'autres contributions, qui
se sont étalées sur les deux semestres. Elles n'ont pas le volume des chapitres précédents, mais
elles occupent une part réelle de mon temps et elles éclairent une autre facette du poste, celle
du maintien en conditions opérationnelles de ce qui existe déjà. J'en présente ici deux : le
maintien du système décisionnel et la contribution au déploiement du SSO.

\section{Le maintien du système décisionnel}

Le groupe pilote son activité à travers MyReport, l'outil décisionnel présenté au chapitre~1.
Chaque service y dispose de tableaux de bord alimentés par des extractions des bases sources.
Comme il n'existe pas d'équipe data constituée, c'est à moi qu'ont été confiés leur maintien et
leur évolution, pour l'ensemble des services demandeurs.

Ce travail prend deux formes. La première est la correction : un tableau ne remonte plus, un
chiffre paraît faux à un utilisateur, une extraction tombe en erreur après une évolution de la
base source. Il faut alors remonter la chaîne de traitement pour identifier l'origine du
problème avant de le corriger, et cette investigation prend souvent plus de temps que la
correction elle-même. La seconde est l'évolution : un service demande un indicateur
supplémentaire, un nouvel axe d'analyse ou une règle de calcul différente.

Ces demandes ont ceci de particulier qu'elles arrivent rarement sous forme de spécification. Un
service décrit une gêne, et c'est à moi de la traduire en règle calculable. C'est le point sur
lequel j'ai le plus progressé sur cette activité, et il rejoint directement ce que j'ai vécu sur
les deux missions principales.

\subsection{Un exemple : la règle de facturation des commandes}

L'exemple le plus représentatif de cette traduction m'a été demandé par le service juridique. La
règle métier est simple à énoncer : pour un lot donné, les commandes se facturent en alternance,
la première étant payante, la deuxième gratuite, la troisième payante, et ainsi de suite.

Elle est cependant beaucoup moins simple à calculer, pour une raison qui tient à la structure de
la donnée. La base ne connaît que des lignes de commande, sans position dans une séquence : rien
n'y indique qu'une commande est la première ou la deuxième de son lot. La catégorie sur laquelle
repose toute la règle, le rang de la commande, n'existe tout simplement pas dans les données. Il
fallait donc la construire.

À cette difficulté s'ajoutait une contrainte pratique. Je ne disposais pas de droits d'écriture
sur la base, et je ne pouvais pas non plus ajouter une colonne aux fichiers sources, qui sont
alimentés par ailleurs. La règle devait donc être portée entièrement par le calcul, sans rien
stocker.

La solution retenue, que représente la figure~\ref{fig:facturation}, tient en une chaîne de
quatre étapes : partir des commandes du lot telles qu'elles sont dans la base, les trier
chronologiquement par opération et par lot, leur attribuer un rang au calcul, puis en déduire le
statut, les rangs impairs étant payants et les rangs pairs gratuits. Aucune de ces étapes n'est
stockée : la chaîne est rejouée à chaque génération du tableau, ce qui permet de tenir la règle
sans droit d'écriture ni colonne ajoutée aux fichiers sources.

\begin{figure}[!htbp]
\centering
\includegraphics[width=0.85\textwidth]{fig_4_1_regle_facturation.png}
\caption{La règle de facturation des commandes, reconstituée par le calcul à chaque génération du tableau}
\label{fig:facturation}
\end{figure}

Ce cas m'a marquée parce qu'il illustre en petit ce que j'ai rencontré en grand sur les deux
missions principales. Le service juridique parle de commandes qui se suivent ; la base ne connaît
que des lignes sans ordre. Entre les deux, il faut construire la catégorie manquante, et c'est ce
travail de construction qui décide de ce que l'entreprise pourra observer ensuite.

\section{La contribution au déploiement du SSO}

Ma seconde activité complémentaire relève davantage de la gestion d'applications que de la
donnée. Le groupe a engagé le déploiement du SSO (Single Sign-On), c'est-à-dire de
l'authentification unique, qui permet à un collaborateur d'accéder à plusieurs applications avec
un seul compte, sans avoir à s'identifier à chaque fois.

J'ai contribué à ce déploiement sur deux périmètres : l'agence de Perpignan et le service
financier. Concrètement, mon rôle a consisté à préparer les postes, à accompagner les
utilisateurs lors de la bascule et à traiter les cas qui ne passaient pas, chaque poste ayant
son historique de comptes et d'accès.

Cette activité peut sembler éloignée du reste de mon mémoire, et elle l'est sur le plan
technique. Elle m'a toutefois appris deux choses que je n'aurais pas apprises ailleurs. La
première est qu'un déploiement se juge à l'usage : une bascule techniquement réussie qui laisse
trois personnes incapables de se connecter est un échec pour le service concerné. La seconde est
qu'une part importante du travail consiste à accompagner, à expliquer et à rassurer, ce qui est
exactement ce que j'ai retrouvé lors de la mise en service du Copilote Financier.

\section{Ce que ces activités apportent au reste}

Ces deux activités partagent un point commun avec mes missions principales : elles portent sur
l'existant, sur ce qui est déjà utilisé au quotidien par des collègues. Une erreur y a donc un
effet immédiat, ce qui impose une prudence constante et une vérification systématique avant toute
mise en service.

Elles m'ont aussi donné une visibilité sur l'ensemble des services du groupe que mes deux
missions principales, plus resserrées sur la direction financière et sur les équipes
commerciales, ne m'auraient pas donnée. C'est en traitant ces demandes que j'ai compris comment
chaque service utilise réellement ses données, et cette connaissance m'a directement servi au
moment de cadrer le Copilote Financier.
